\section{Heliocentric checks and worked bounded cases}
\label{sec:helio}
Use $R_0=149597870700$ m and
$\mu_\odot=1.32712440018\,10^{20}$ $\mathrm{m^3/s^2}$.
Then $T_0=58.13244087$ days, $V_0=29.78469183$ km/s and
$J_0=176.6257101$ $\mathrm{m^2/s^3}$.

\subsection{Revisiting the Jupiter-class candidates of the transfer article}
The ideal engine of \cite{OriginalKepler} has $P=1$ GW, $m_0=3\,10^6$ kg and
$m_d=10^6$ kg: $\chi=3.774459936$ and $\kappa=7.548919872$.
The public atlas of \cite{OriginalKepler} supplies warm starts, not
certified trajectories. We independently integrate the normalized Cartesian state,
effort adjoints and effort quadrature, and solve four endpoint equations
plus $J=\kappa$. This avoids relying on the repository viewer's radius
and angle acceptance tolerances. The geometries below reproduce the
short-arc table and the two illustrated solar-pass families of the
transfer article, allowing for its rounded figure labels.

\begin{center}\small
\begin{tabular}{rrrrr}\toprule
$\rho,\Theta$&days&$r_{\min}$ AU&$v_{e,\min}$&$v_{e,\max}$\\\midrule
$5.2,80^\circ$&170.034&1.0000&40.60&2984.8\\
$5.0,273.1^\circ$&206.560&0.16993&42.35&13001.0\\
$5.2,385.2^\circ$&214.236&0.000273&4.85&1203.5\\\bottomrule
\end{tabular}\end{center}
Exhaust speeds are in km/s. All three paths fail the band. The first two
have numerical radius minima above 0.1 AU; the last also fails the floor.
Scaled endpoint residuals are below $4\,10^{-13}$ for the first two and
$10^{-8}$ for the last. The long-arc solve uses polar equations and
analytic shooting sensitivities to preserve its winding; a radial-velocity
root refines its narrow perihelion, which uniform samples alone miss.


We also reconstruct all 73 rows of that article's $5.2$ AU final-polar-angle
table, $\Theta=0^\circ,5^\circ,\ldots,360^\circ$, by independent continuation.
Every row violates the exhaust band. Their reconstructed times range from
$169.7643$ to $220.4991$ days (the $0^\circ$ row lies outside the abbreviated
170--215 day range). Scaled endpoint residuals are below $4.8\,10^{-11}$.
The radial audit locates turning points rather than only testing a time grid:
\begin{center}\small
\begin{tabular}{lrr}\toprule
Final-angle rows&Band rejected&Floor rejected\\\midrule
$0^\circ$--$300^\circ$ (61 rows)&61&0\\
$305^\circ$--$360^\circ$ (12 rows)&12&12\\\bottomrule
\end{tabular}\end{center}
For example, the numerical minimum radius is $0.1011887$ AU at $300^\circ$
and $0.090577$ AU at $305^\circ$. Thus none of these 73 candidates, nor
either illustrated solar-pass candidate, survives both operational checks.

The audit evaluates $v_e=V_0(\chi+I)/|u|$ and the radius on a dense
40001-point reconstruction, with continuous refinement of the long-arc
perihelion and exhaust extrema. Each reported band rejection has a large
violation on a continuous powered interval, so it does not depend on
interpreting an isolated zero-thrust sample. Clipping a failed trajectory
to the band does not repair its endpoint. The input seeds and normalized initial covectors are provided as
supplementary material.

\subsection{A local certificate demonstration: a 2.9-day solar hop}
This example raises the heliocentric radius from $1$ to $1.001$ AU
in $2.90663$ days and ends on a slightly eccentric orbit. Its purpose
is to demonstrate an explicit feasible path together with global effort
and time bounds on a small, fully checkable case.
For this auditable local example, keep $P=1$ GW and $m_0=3\,10^6$ kg,
but choose the smaller resource $\kappa=0.0320066666667$. Begin on the
unit circular orbit and let $c=0.4$, $H=0.05$. Define a physical
heliocentric path through its normalized polar coordinates:
\begin{equation}
 \begin{gathered}
 r(t)=1+ct^2,\qquad \theta(t)=\frac{t}{\sqrt{1+ct^2}},\\
 v_r=2ct,\qquad v_\theta=r^{-1/2}.
 \end{gathered}
\end{equation}
Since $\theta'=r^{-3/2}$, exact substitution into the Kepler equations
gives $u_r=2c$ and $u_\theta=ct/r^{3/2}$. Thus this is a continuously
powered exact path with $r\ge1$ AU, not a patched approximation. Its
endpoint defines the target: take the unforced Kepler orbit through
$(r(H),\theta(H),v_r(H),v_\theta(H))$ at epoch $H$. The arrival orbit is
slightly eccentric.
Coasting along that orbit is admissible on the short horizon used below.
The same backward comparison used below bounds its separation from the
unit circle by $0.0031$ on $[0,H]$, so the target orbit also clears the floor.

The effort and a convenient elementary upper bound are
\begin{equation}
\begin{split}
 J&=4c^2H+c^2\int_0^H\frac{t^2}{(1+ct^2)^3}dt\\
  &=0.0320066546838,\\
 U&=4c^2H+c^2H^3/3=0.0320066666667=\kappa.
\end{split}
\end{equation}
Mass accounting gives $m_d=2974774.5061$ kg from the budget and an actual
propellant expenditure $25225.4845$ kg for this witness. The burn satisfies
$0.8\le|u|\le\sqrt{0.64+0.16H^2}$ and $\chi\le w\le\chi+U$.
Consequently its full-power exhaust remains between $140.4825$ and
$141.7181$ km/s, strictly within the required band, throughout the path.

Here is a global, conservative lower-bound calculation, which also
brackets the minimum time instead of assigning optimality to this shaped
path. Compare every competitor of effort at most $U$ with the unit
circular coast. In a position tube of radius $d=0.1$, the gravity
Lipschitz constant is $K=2/0.9^3$, $\omega=\sqrt K$.
Variation of constants and Cauchy--Schwarz give
\begin{equation}
 |\delta r(t)|\le k_r(t)\sqrt J,\qquad
 |\delta v(t)|\le k_v(t)\sqrt J,
\end{equation}
where $k_r,k_v$ are those in Section~\ref{sec:verification}. Since
$\sqrt U k_r(H)<0.0012<0.1$, the first-exit argument validates this tube
for every cheaper competitor, irrespective of its thrust direction.
The target endpoint has $|\delta v(H)|>0.0399$ and
$k_v(H)^2\le H\cosh^2(\omega H)<1.007H$. Therefore
\begin{equation}
 0.0316=L\le F^*(0,H)\le U=0.0320066666667.
\end{equation}
The clipping argument from \eqref{eq:globalL} covers more expensive
competitors as well. These same bounds hold for the restricted engine,
because its witness passes the band and its competitors are a subset of
the ideal ones.

For an earlier arrival at $h_-=0.049$, compare the ballistic target to
the unit coast backward from $H$. At $H$, $|\delta r|<0.00101$ and
$|\delta v|<0.041$. Over $H-h_-=0.001$, the comparison inequalities
give $|\delta r|<0.0011$ and a change in $\delta v$ of less than
$K(0.0011)(0.001)<0.0000031$. Hence every affordable earlier competitor
would need
\begin{equation}
 J\ge\frac{(0.0399-0.0000031)^2}{1.007(0.049)}
   >0.0322591>\kappa,
\end{equation}
a contradiction. Coast extension and the monotone value then exclude
all still earlier arrivals. Thus the resource-to-time-to-fuel calculation is
\begin{equation}
 \boxed{\begin{gathered}0.049<h^*(\kappa)\le0.05,\\
 2.84848<T^*\le2.90663\ {\rm days}.\end{gathered}}
\end{equation}
The stated fuel is that of the witness path; the optimal duration lies
somewhere in this interval. The conservative decimal
inequalities can be checked by rational Taylor enclosures; a supplementary
checker does so for the bounds used in the time bracket.


\subsection{A certified 70-day Earth--Mars witness}
\label{sec:marswitness}
For an interplanetary performance point, return to $P=1$ GW,
$m_0=3\,10^6$ kg and $m_d=10^6$ kg, hence
$\kappa=7.548919872$. Use coplanar circular planetary orbits at
$1$ and $R_M=1.524$ AU in the solar two-body model. The following
explicit path gives full position-and-velocity rendezvous in exactly
70 days, with useful jet power fixed at 1 GW throughout. The exhaust
schedule is variable and satisfies the two-sided $50$--$250$ km/s band.

Let $H=70(86400)/T_0$ and $q=0.51011006$, with the decimal exact.
Set $c_q=(1-q^2)/(1+q^2)$ and $s_q=2q/(1+q^2)$. The endpoint states are
\begin{equation}
\begin{gathered}
 r_0=(1,0),\qquad v_0=(0,1),\\
 r_8=R_M(c_q,s_q),\qquad
 v_8=R_M^{-1/2}(-s_q,c_q).
\end{gathered}
\end{equation}
The circular target passes through this endpoint at $H$. Thus its
initial phase is $2\arctan q-H/R_M^{3/2}$, approximately
$17.3820101^\circ$ ahead of Earth; arrival longitude is approximately
$54.0531714^\circ$. Planetary
escape, capture, parking-orbit costs and departure waiting time are
outside this heliocentric model.

\begin{table*}[t]
\centering\small
\caption{Exact rational interior knots of the 70-day witness. Entries
are exact defining decimals. Positions and velocities use the solar normalized units.
Endpoint knots are $t_0=0$, $t_8=H$ and the analytic states in the text.}
\label{tab:marsknots}
\begin{tabular}{rrrrrr}\toprule
$j$ & $t_j/H$ & $r_{jx}$ & $r_{jy}$ & $v_{jx}$ & $v_{jy}$\\\midrule
1 & 0.137169299389 & 1.001350238628 & 0.168727459300 & 0.017525621592 & 1.046344765303\\
2 & 0.274338598779 & 1.006411818893 & 0.346210491047 & 0.046245191438 & 1.105077859977\\
3 & 0.411507898168 & 1.017566072445 & 0.534820743099 & 0.091838207401 & 1.182668462417\\
4 & 0.548677197557 & 1.037723200329 & 0.738442255030 & 0.154970973170 & 1.288219089473\\
5 & 0.661507898168 & 1.043707872696 & 0.902107682569 & -0.064178195495 & 1.117609239693\\
6 & 0.774338598779 & 1.020891188142 & 1.041194789559 & -0.269833623616 & 0.926300052106\\
7 & 0.887169299389 & 0.970819319098 & 1.152810168550 & -0.465866897104 & 0.712855383845\\
\bottomrule\end{tabular}
\end{table*}

On interval $j=0,\ldots,7$, put $\Delta_j=t_{j+1}-t_j$,
$s=(t-t_j)/\Delta_j$ and use cubic Hermite interpolation:
\begin{equation}
\begin{split}
 r(t)={}&(2s^3-3s^2+1)r_j
       +(s^3-2s^2+s)\Delta_j v_j\\
 &+(-2s^3+3s^2)r_{j+1}
       +(s^3-s^2)\Delta_j v_{j+1}.
\end{split}
\label{eq:marshermite}
\end{equation}
Table~\ref{tab:marsknots} and the analytic endpoints specify the whole
path. Define its control and mass by
\begin{equation}
\begin{gathered}
 u=r''+r/|r|^3,\qquad I(t)=\int_0^t|u|^2\,ds,\\
 m(t)=\frac{m_0\chi}{\chi+I(t)},\qquad
 v_e(t)=V_0\frac{\chi+I(t)}{|u(t)|}.
\end{gathered}
\end{equation}
The position and velocity are continuous and the endpoint conditions
hold exactly by construction. Acceleration may jump at knots; the
model imposes no thrust-slew constraint. Each one-sided control value
passes the band. The resulting mass, thrust and exhaust obey the
constant-power equations exactly almost everywhere.

A whole-interval enclosure gives
\begin{equation}
\begin{gathered}
 2.5212<J<2.5307< U_M=2.6<\kappa,\\
 99<v_e(t)<102\ {\rm km/s},\qquad |r(t)|>0.999\ {\rm AU},\\
 1.2014\,10^6<m_0-m(H)<1.2041\,10^6\ {\rm kg}.
\end{gathered}
\end{equation}
In particular, the spacecraft retains over $1.7959\,10^6$ kg, above
the dry-mass floor. Consequently, for this target phase,
\begin{equation}
 \boxed{F_{\rm band}^*(0,H)\le U_M=2.6,\qquad T_{\rm band}^*(\kappa)\le70\ {\rm days}.}
\end{equation}
This is an unoptimized upper bound on the minimum time. Comparing it
with an ideal-engine optimum for the same endpoints, to measure the
time cost of the band, is left for future work.

Here is a reproducible enclosure prescription. Form the Hermite cubic
coefficients from the defining decimals, using outward interval arithmetic
at 60 significant decimal digits, including the square roots in $T_0$,
$V_0$ and the terminal velocity. Subdivide each of the eight local
$s$ intervals into 128 equal parts. Horner evaluation encloses $r$ and
$r''$; square, sum and take an outward square root to bound $|r|$,
then enclose $u$ and $|u|^2$. Multiplying each latter enclosure by its
time width gives lower and upper Riemann bounds for $I$. For the
exhaust check on a subinterval, use the cumulative lower bound at its
start and cumulative upper bound at its end. This monotonic enclosure
covers every time within it. The resulting bounds imply the displayed
coarse inequalities. A standalone supplementary checker implements
this recipe; a separate floating-point quadrature and dynamics
integration provide independent numerical cross-checks.
